% =============================================================================
% ISI Country Brief — France
% External Supplier Concentration Profile
%
% ISI Paper Series — Country Brief
% International Sovereignty Institute — Research Unit
% March 2026
% =============================================================================
\documentclass[12pt,a4paper]{article}

% --- ISI Paper Series shared template ----------------------------------------
\usepackage{isi-series}

% --- PDF metadata ------------------------------------------------------------
\hypersetup{
  pdftitle={External Supplier Concentration Profile: France — ISI EU-27 2024},
  pdfauthor={International Sovereignty Institute -- Research Unit},
  pdfcreator={International Sovereignty Institute},
  pdfsubject={ISI Paper Series -- Country Brief FR, EU-27, Vintage 2024},
  pdfkeywords={sovereignty index, EU-27, supplier concentration,
    france, country brief},
}

% ==============================================================================
\begin{document}
\hypersetup{pageanchor=false}

% ---------- Title Page with Metadata ------------------------------------------
\begin{titlepage}
\centering
\vspace*{2cm}
{\Large\textsc{ISI Paper Series --- Country Brief}}\\[1.2cm]
{\LARGE\bfseries External Supplier Concentration Profile:\\[0.3em]
France}\\[1.5cm]
{\large International Sovereignty Index (ISI)\\[0.3em]
EU-27 --- Vintage 2024 --- Methodology v1.0}\\[1.2cm]
{\large International Sovereignty Institute --- Research Unit}\\[0.8cm]
{\large March 2026}

\vspace{0.8cm}
\noindent\rule{\textwidth}{0.4pt}
\vspace{0.4cm}

% --- Research Metadata --------------------------------------------------------
\begin{flushleft}
\noindent\textbf{Data basis:} ISI Paper~2 --- EU-27 Empirical Results 2024 (Methodology~v1.0, Vintage~2024).\\[4pt]
\noindent\textbf{Methodology:} ISI Paper~1 --- Methodological Foundations, Version~1.0.

\medskip

\noindent\textbf{JEL Codes:} F02, F15, F52, O30, Q43\\[4pt]
\noindent\textbf{Keywords:} sovereignty index; EU-27; supplier concentration; composite indicator; variance decomposition; defence dependence

\medskip

\noindent\textbf{Related publications in the ISI Paper Series:}
\begin{itemize}[nosep,leftmargin=1.5em]
\item ISI Paper~1 --- Technical Specification (v1.0) --- doi:\\
  \url{https://doi.org/10.5281/zenodo.18764227}
\item ISI Paper~2 --- EU-27 Empirical Results 2024 (v1.0) --- doi:\\
  \url{https://doi.org/10.5281/zenodo.18764170}
\end{itemize}
\end{flushleft}
\vfill
\end{titlepage}

\hypersetup{pageanchor=true}
\pagenumbering{arabic}
\pagestyle{fancy}
\fancyhead[L]{\small\sffamily\textit{ISI Country Brief}}
\fancyhead[R]{\small\sffamily\textit{France}}
\renewcommand{\headrulewidth}{0.4pt}

% ---------- Body --------------------------------------------------------------

\section{Executive Summary}
\label{sec:summary}

France ranks~27th in the EU-27 with a composite \ISI score of~0.236, classified as \emph{mildly concentrated}.  The dominant axes are defence~(0.370) and logistics~(0.353).  The structural profile is characterised as dual-axis concentration (defence and logistics).  Under Dirichlet weight perturbation (10\,000~Monte Carlo draws), France's mean rank is~26.55 with a standard deviation of~0.50; the 5th--95th percentile band spans ranks~26 to~27.


\section{Country Position in the EU-27}
\label{sec:position}

France occupies rank~27 of~27 EU member states (composite score~0.236).  The EU-27 mean composite score is~0.344 with a standard deviation of~0.070.  France's score lies below the EU-27 mean, indicating below-average external supplier concentration.  France holds the lowest rank in the EU-27.  The next country above is Slovakia (rank~26; 0.236).


\section{Axis Profile}
\label{sec:axis-profile}

\Cref{tab:axis-profile} presents France's scores on all six \ISI axes.

\begin{table}[htbp]
\centering\small
\caption{France's \ISI axis profile.}
\label{tab:axis-profile}
\begin{tabular}{@{}lr@{}}
\toprule
\textbf{Axis} & \textbf{Score} \\
\midrule
    Finance           & 0.100 \\
    Energy            & 0.309 \\
    Technology        & 0.120 \\
    Defence           & 0.370 \\
    Critical Inputs   & 0.161 \\
    Logistics         & 0.353 \\
\bottomrule
\end{tabular}
\end{table}

The highest axis score is defence~(0.370), which lies below the EU-27 axis mean of~0.573.  The second-highest axis is logistics~(0.353), below the EU-27 axis mean of~0.518.  The lowest axis score is finance~(0.100), below the EU-27 mean of~0.148.


\section{Structural Drivers}
\label{sec:drivers}

\begin{sloppypar}
France's composite score is driven by two axes in combination: defence~(0.370) and logistics~(0.353).  The gap between these two axes is~0.017, indicating a dual-axis concentration profile.  Neither axis alone dominates the composite; rather, the combination of elevated scores on both dimensions determines France's position in the ranking.
\end{sloppypar}


\section{Comparative Context}
\label{sec:comparative}

Across the EU-27, the covariance-based variance decomposition shows that defence (35.1\,\%) and logistics (34.3\,\%) together account for 69.4\,\% of total cross-sectional composite variance.  France's defence score~(0.370) lies below the EU-27 defence mean~(0.573), and its logistics score~(0.353) lies below the EU-27 logistics mean~(0.518).  Under Dirichlet weight perturbation, France's rank standard deviation is~0.50, indicating high rank stability.



France records the lowest composite score in the EU-27, a position that reflects broad-based supplier diversification across all six axes.  Finance~(0.100) and technology~(0.120) are the two lowest axis scores, both well below the EU-27 means.  Defence~(0.370) is the highest axis in the profile but still falls below the EU-27 defence mean of~0.573.  The absence of any axis exceeding the EU-27 average produces the lowest composite in the dataset.

The rank standard deviation of~0.50 and the one-position percentile band~(26--27) make France's rank among the most stable in the EU-27.  Under no Dirichlet weight configuration does France rise above rank~26.  This extreme stability is a consequence of consistent below-mean positioning on every axis: regardless of how weights are distributed, the composite remains in the lowest tier.

The co-occurrence of France and Slovakia at the bottom of the distribution---both with composites of~0.236---should not obscure their structural differences.  Slovakia's low composite is driven by a single zero-valued axis, whereas France's reflects uniformly low concentration across all dimensions.  The distinction is analytically important: the two countries arrive at the same composite score through fundamentally different axis structures.


\section{Interpretation Boundary}
\label{sec:interpretation}

The \ISI measures concentration of external suppliers, not sovereignty in a normative or political sense.  High concentration may reflect geography, economic specialisation, or alliance structures.  A high composite score does not imply policy failure; a low score does not imply policy success.  The index provides an empirical measurement basis --- what policymakers derive from it is a separate question.


% ---------- References --------------------------------------------------------
\clearpage
\vspace*{1.5cm}

\begingroup\emergencystretch=2em\hbadness=10000
\section*{References}

\begin{enumerate}[label={[\arabic*]},leftmargin=2em,itemsep=6pt]
\item International Sovereignty Institute.
  \textit{International Sovereignty Index (ISI) technical specification, version~1.0}.
  Zenodo, 2026.\\ISI Paper Series, Paper~1.
  doi:\,\href{https://doi.org/10.5281/zenodo.18764227}{10.5281/zenodo.18764227}.

\item International Sovereignty Institute.
  \textit{International Sovereignty Index (ISI): EU-27 empirical results 2024 (v1.0)}.
  Zenodo, 2026.\\ISI Paper Series, Paper~2.
  doi:\,\href{https://doi.org/10.5281/zenodo.18764170}{10.5281/zenodo.18764170}.
\end{enumerate}
\endgroup

\end{document}
